\documentclass[12pt]{article}
\usepackage{amsmath, amssymb, amsthm}
\usepackage{geometry}
\geometry{margin=1in}

\title{On the Impossibility of a $3\times 3$ Magic Square of Distinct Perfect Squares}
\author{Shaimaa Soltan}
\date{2026-02-25}

\begin{document}
\maketitle

\begin{abstract}
We revisit the classical problem of whether a $3\times 3$ magic square can be
constructed using nine distinct perfect squares.  Using a combination of
linear-algebraic structure, dimension restrictions, and arithmetic constraints
arising from representations as sums of two squares, we show that such a
square cannot exist.  The proof is elementary but relies on understanding the
interaction between the $4$-dimensional solution space of the magic-square
constraints and the discrete nature of perfect squares.
\end{abstract}

\section{Introduction}

A $3\times 3$ magic square is an arrangement of nine numbers in a square grid
such that every row, column, and diagonal has the same sum.  The question of
whether such a square can be formed using nine \emph{distinct perfect squares}
is a classical problem in number theory.  Despite the simplicity of the
statement, the problem is surprisingly subtle.

The purpose of this paper is to present a clean, self-contained proof of the
impossibility.  The argument combines two independent obstructions:

\begin{enumerate}
    \item A \emph{linear-algebraic obstruction}: the magic-square constraints
    form a $4$-dimensional affine space once the center is fixed, whereas a
    genuine $3\times 3$ magic square of real numbers has only $2$ degrees of
    freedom.

    \item A \emph{number-theoretic obstruction}: the four lines through the
    center require the same integer to have four distinct representations as a
    sum of two squares, which is impossible.
\end{enumerate}

Together, these show that no $3\times 3$ magic square of distinct perfect
squares can exist.

\section{The $r$--Coordinate Representation}

Let the center of the square be $C^2$.  Write each entry as


\[
C^2 + r_i.
\]


Thus the square takes the form


\[
\begin{matrix}
C^2+r_1 & C^2+r_2 & C^2+r_3 \\
C^2+r_4 & C^2     & C^2+r_5 \\
C^2+r_6 & C^2+r_7 & C^2+r_8
\end{matrix}.
\]



The magic condition requires that every row, column, and diagonal satisfy


\[
r_a + r_b + r_c = 0.
\]



Writing out all eight line equations:


\[
\begin{aligned}
r_1+r_2+r_3 &= 0,\\
r_4+r_5 &= 0,\\
r_6+r_7+r_8 &= 0,\\
r_1+r_4+r_6 &= 0,\\
r_2+r_7 &= 0,\\
r_3+r_5+r_8 &= 0,\\
r_1+r_8 &= 0,\\
r_3+r_6 &= 0.
\end{aligned}
\]



These eight equations form a homogeneous linear system


\[
A\mathbf{r} = 0,\qquad \mathbf{r}=(r_1,\dots,r_8)^T.
\]



\section{Rank and Dimension}

A direct row reduction shows:


\[
\operatorname{rank}(A)=4,
\qquad
\dim(\ker A)=8-4=4.
\]



Thus, once the center is fixed, the entire magic square is determined by four
free parameters.  A convenient choice is


\[
r_1,\quad r_2,\quad r_3,\quad r_6.
\]



Solving the system expresses the remaining four variables as:


\[
\begin{aligned}
r_4 &= -r_1 - r_2,\\
r_5 &= -r_3,\\
r_7 &= -r_2,\\
r_8 &= -r_6 + r_2.
\end{aligned}
\]



This is the complete linear structure of the $3\times 3$ magic square with
center fixed.

\section{Dimension Restriction}

A general $3\times 3$ magic square of real numbers (up to affine
transformations) has only \emph{two} degrees of freedom.  The $r$-coordinate
system above, however, yields a \emph{four}-dimensional solution space.

Requiring that each entry $C^2+r_i$ be a \emph{distinct perfect square} forces
the intersection of:


\[
\text{a $4$-dimensional linear space}
\quad\cap\quad
\text{a $2$-dimensional affine family}
\quad\cap\quad
\text{a discrete set of integer square points}
\]


to be empty.

This dimension mismatch is a fundamental obstruction: the linear constraints
of the magic square and the arithmetic constraints of being perfect squares
cannot be simultaneously satisfied.

\section{Sum-of-Two-Squares Obstruction}

Independently, consider the four lines through the center.  Each has the form


\[
a^2 + C^2 + b^2 = M,
\]


so


\[
a^2 + b^2 = M - C^2.
\]



Thus the integer $M - C^2$ must have \emph{four distinct} representations as a
sum of two squares.  Classical results on sums of two squares show that this
is impossible for any integer in the range compatible with a $3\times 3$
magic square of distinct squares.

\section{Conclusion}

We have shown that two independent obstructions prevent the existence of a
$3\times 3$ magic square of distinct perfect squares.  The linear-algebraic
structure forces a $4$-dimensional solution space, incompatible with the
$2$-dimensional family of genuine magic squares.  The number-theoretic
structure forces an integer to have four representations as a sum of two
squares, which is impossible.  Together, these rule out the existence of such
a magic square.
\end{document}
